Let $\Omega$ be a set of linear transformations on a vector space $\mathfrak{R}$, and let $\mathfrak{S}$ be an invariant subspace of $\mathfrak{R}$ relative to $\Omega$. Let $\mathfrak{U}$ be an invariant subspace of $\mathfrak{R}$ containing $\mathfrak{S}$.

Then $\mathfrak{U}$ is said to be **irreducible over $\mathfrak{S}$** (relative to $\Omega$) if the factor space $\mathfrak{U}/\mathfrak{S}$, together with the set $\bar{\Omega}$ of linear transformations induced by $\Omega$, admits no proper nonzero invariant subspace. Equivalently, there is no invariant subspace $\mathfrak{W}$ such that $\mathfrak{U} \supsetneq \mathfrak{W} \supsetneq \mathfrak{S}$.

In terms of matrices: if $(\beta)$ is the collection of matrices representing the induced transformations on $\mathfrak{U}/\mathfrak{S}$ relative to some basis, then irreducibility means that it is impossible to find an invertible matrix $(\mu)$ such that all the conjugates $(\mu)(\beta)(\mu)^{-1}$ simultaneously take the same nontrivial reduced block form.
